\section{Speicherung und Weitergabe von Informationen im historischen Kontext}
\label{sec:einleitung}

Informationen speichern und weitergeben ist zentraler Bestandteil der
menschlichen Kultur und führte in der Menschheitsgeschichte zu bahnbrechenden
Erfindungen und laufenden Innovationen.

In der Steinzeit wurden Informationen an Höhlenwände gemalt, ab dem
\nolbreaks{3. Jahrtausend} vor Christus mittels Keilschrift (siehe
\autoref{fig:keilschrift}) in Steintafeln gemeißelt und in der Antike auf
Papyrus gezeichnet. Mitte des \nolbreaks{15. Jahrhunderts} revolutionierte
Gutenberg die Informationsweitergabe mit der Erfindung des Buchdrucks, der die
Produktion von Büchern in hohen Stückzahlen und zu geringen Kosten ermöglichte.
Auch heute noch wird diese Technik für die Reproduktion von Texten verwendet.

Ende des \nolbreaks{19. Jahrhunderts} gelang es erstmals, Musik zu speichern und
abzuspielen. Hierfür wurden zunächst Musikwalzen eingesetzt. Diese wurden um
1900 dann von Schallplatten aus einer Pressmasse aus Schellack\footnote{harzige
Ausscheidung von Lackschildläusen \cite{schellack}} (siehe
\autoref{fig:schallplatte-schellack}) abgelöst. Auf Schellackplatten konnten
Musikstücke mit bis zu fünf Minuten Spielzeit gespeichert werden und wurden auf
einem Grammophon abgespielt. Mit dem Einsatz des Kunststoffes Polyvinylchlorid
ab den 50er Jahren des \nolbreaks{20. Jahrhunderts} verbesserte sich die
Tonqualität von Schallplatten deutlich. Außerdem konnten längere Stücke
gespeichert und wiedergegeben werden. \cite{schallplatte1}

Mit der Compact Cassette (CC bzw. MC, siehe \autoref{fig:compact-cassette}), die
1963 auf den Markt kam, war es dann jedem möglich, selbst Musik aufzunehmen und
dauerhaft zu speichern \cite{kassette}. Compact Cassetten verwenden ein
Magnetband, das aus einer langen schmalen Kunststofffolie besteht, die mit einem
magnetisierbaren Material beschichtet wurde \cite{tonband}.

\ifthenelse{\boolean{showPics}}{
    \begin{figure}[h]
        \begin{center}
            \begin{minipage}[t]{0.3\textwidth}
                \begin{center}
                    \includegraphics[height=0.1\textheight]{Bilder/Einleitung/keilschrift.png}
                    \caption[Altpersische Keilschrift \newline \url{http://www.faz.net/aktuell/feuilleton/geisteswissenschaften/keilschriften-zehntausend-freunde-mesopotamiens-1657070.html} (zuletzt aufgerufen am 19.09.2015)]{Altpersische Keilschrift}
                    \label{fig:keilschrift}
                \end{center}
            \end{minipage}
            \hspace{0.025\textwidth}
            \begin{minipage}[t]{0.3\textwidth}
                \begin{center}
                    \includegraphics[height=0.1\textheight]{Bilder/Einleitung/schallplatte-schellack.png}
                    \caption[Schellackschallplatte \newline \url{https://de.wikipedia.org/wiki/Datei:Schallplatte_Deutsche_Grammophon_Stimme_seines_Herrn.jpg} (zuletzt aufgerufen am 19.09.2015)]{Schellackschall-platte}
                    \label{fig:schallplatte-schellack}
                \end{center}
            \end{minipage}
            \hspace{0.025\textwidth}
            \begin{minipage}[t]{0.3\textwidth}
                \begin{center}
                    \includegraphics[height=0.1\textheight]{Bilder/Einleitung/compact-cassette.png}
                    \caption[Compact Cassette \newline \url{https://de.wikipedia.org/wiki/Datei:Compactcassette.jpg} (zuletzt aufgerufen am 19.09.2015)]{Compact Cassette}
                    \label{fig:compact-cassette}
                \end{center}
            \end{minipage}
        \end{center}
    \end{figure}
}{}

In den 80er Jahren kam die Compact Disc (CD, siehe \autoref{fig:compact-disc})
auf den Markt. Hierbei handelt es sich um einen optischen Datenträger aus dem
Kunststoff Polycarbonat handelt. Compact Discs und Digital Video Discs (DVD),
eine Weiterentwicklung der CD, zeichnen sich durch hohe Speicherkapazität sowie
geringe Abnutzungserscheinungen aus. Deshalb wurden sie zum universellen
Speichermedium für Musik, Dokumente, Bilder und Filme.

Mit der Verbreitung des Internets hat die Übertragung von Informationen mittels
Lichtsignalen in den letzten Jahren an Bedeutung gewonnen. Neben Glasfaserkabeln
werden hierfür polymer optische Fasern (POF, siehe \autoref{fig:pof})
eingesetzt.

Anhand der Beispiele der Compact Disc und der polymer optischen Faser wird in
dieser Arbeit der Einsatz von transparenten Kunststoffen in der
Informationstechnologie aufgezeigt und ihre spezielle Eignung für die
Informationsspeicherung und -weitergabe erörtert. Dabei wird auf die
Funktionsweise und Produktion von optischen Datenträgern und optischen
Wellenleitern eingegangen. Außerdem werden die physikalischen Eigenschaften und
die Herstellung der verwendeten Kunststoffe erläutert.

\ifthenelse{\boolean{showPics}}{
    \begin{figure}[h]
        \begin{center}
            \begin{minipage}[t]{0.4\textwidth}
                \begin{center}
                    \includegraphics[height=0.1\textheight]{Bilder/Einleitung/compact-disk.png}
                    \caption[Compact Disc \newline \url{https://en.wikipedia.org/wiki/File:Compact_disc.svg} (zuletzt aufgerufen am 19.09.2015)]{Compact Disc}
                    \label{fig:compact-disc}
                \end{center}
            \end{minipage}
            \hspace{0.025\textwidth}
            \begin{minipage}[t]{0.4\textwidth}
                \begin{center}
                    \includegraphics[height=0.1\textheight]{Bilder/Einleitung/pof.png}
                    \caption[polymer optische Faser \newline \url{http://www.heise.de/tr/imgs/08/2/5/4/2/1/5/c337ef89957e0f2b.jpg} (zuletzt aufgerufen am 19.09.2015)]{polymer optische Faser}
                    \label{fig:pof}
                \end{center}
            \end{minipage}
        \end{center}
    \end{figure}
}{}
